\sgaStableTarget{sga6.e12.sec.1}\section*{1. Statement of the main theorem and applications}

\sgaStableTarget{sga6.e12.thm.1-1}\textbf{Theorem 1.1.} \uline{Let} @@M0@@ \uline{be a noetherian scheme,} @@M1@@ \uline{and} @@M2@@ \uline{two proper} @@M3@@\uline{-schemes, and} @@M4@@ \uline{a morphism. Suppose that} @@M5@@ is integral \uline{and} @@M6@@ surjective. \uline{Then there is a nonempty open subset} @@M7@@ \uline{of} @@M8@@ \uline{such that the morphism}
@@M9@@
\uline{between the Picard functors (localized for the} f.p.p.f \uline{topology) is representable by a quasi-affine morphism of finite presentation.}

\sgaStableTarget{sga6.e12.cor.1-2}\textbf{Corollary 1.2.} \uline{Let} @@M10@@ \uline{be a proper scheme over a noetherian scheme} @@M11@@\uline{. Suppose that} @@M12@@ is integral. \uline{Then there is a nonempty open subset} @@M13@@ \uline{of} @@M14@@ \uline{such that the Picard functor} @@M15@@ \uline{is representable by a} @@M16@@\uline{-group scheme that is a sum of quasi-projective open subschemes of finite type, hence a fortiori separated over} @@M17@@.

Indeed, recall the following fact (FGA N\textsuperscript{o}232): let @@M18@@ be a scheme, and let @@M19@@ be a projective @@M20@@-scheme, of finite presentation and integral over @@M21@@ (that is, flat with geometrically integral fibers). Then the Picard functor is represented by an @@M22@@-group scheme @@M23@@, an increasing union of quasi-projective open subschemes of finite presentation.

Consideration of the ``Hilbert polynomials'' makes it possible here to decompose @@M24@@ as a sum of open subschemes @@M25@@, @@M26@@. In the next exposé it is proved, under the additional hypothesis that @@M27@@ is noetherian, that the @@M28@@ are of finite type over @@M29@@. It is then known that 1.2 holds in the case where @@M30@@ is in addition projective and integral over @@M31@@; we shall deduce the general case using the

\sgaStableTarget{sga6.e12.lem.1-3}\textbf{Lemma 1.3.} \uline{Let} @@M32@@ \uline{be of finite type over an integral noetherian} @@M33@@\uline{. Then there are a nonempty open subset} @@M34@@ \uline{of} @@M35@@\uline{, an integral scheme} @@M36@@ \uline{and a finite and}

